\chapter{考察課題の解答の手引き}\label{app:answers}

各章の考察課題について、考え方の道筋と数値の目安を示す。唯一の正解がない課題では、検討すべき観点を挙げる。

\section{基礎編}

\paragraph{剛体力学の基礎 1（胴体の慣性を 2 倍にする）}
$\Gamma=B/A$ で、$A$ の大部分は胴体の慣性 $I_b$（5.6 kg\,m$^2$）が占める。$I_b$ を 2 倍にすると $A$ はおよそ 6.7 → 12.3 kg\,m$^2$ に増え、$B$ は変わらない。
したがって $\Gamma$ はおよそ 0.55 倍になり、腕を 180° 振ったときの回転（16°）も約 9° に減る。

\paragraph{剛体力学の基礎 2（回転の順序）}
$z$ 軸まわり 90° の回転 $R_z$ と $x$ 軸まわり 90° の回転 $R_x$ について、
$R_zR_x$ は $x$ 軸を $y$ 軸へ、$R_xR_z$ は $x$ 軸を $z$ 軸へ移すので、両者は異なる。具体的に行列を書いて確かめるとよい。

\paragraph{幾何力学入門 1（符号の違う曲率をまたぐループ）}
1 周の回転は曲率の面積分なので、正の領域と負の領域の寄与が打ち消し合う。回転を大きくするには、同じ符号の領域だけを囲むループが最適である（付録~\ref{app:connection}）。

\paragraph{幾何力学入門 2（同じ位相の往復）}
腕と脚を同じ位相で往復させると、形状空間では同じ線分を行って戻るだけで、面積を囲まない。線積分は行きと帰りで打ち消し合い、正味の回転はゼロになる。

\paragraph{最適化・制御・学習の基礎 1（周回周波数への小さな費用）}
角運動量を一度入れた後は、AMBAC を使っても使わなくても推進剤は同じになる（第~\ref{ch:ai}章の考察 C-2）。
線形計画の最適解が一意に定まらないので、小さな費用で「不要なループは回さない」解を選ばせる。費用が十分小さければ推進剤の最小値は変わらない。

\paragraph{最適化・制御・学習の基礎 2（ES のノイズの大きさ）}
$\sigma$ が大きすぎると、勾配の推定が「離れた場所の平均」になり、細かい最適化ができず、学習が不安定になる。
小さすぎると、報酬の違いが乱数の揺らぎに埋もれる。特に幾何学的位相は振幅の 2 次の効果なので（第~\ref{ch:basics-geom}章）、静止に近い方策のまわりでは $\sigma$ が小さいと勾配がほとんど見えない。

\section{第 I 部}

\paragraph{第~\ref{ch:theory}章 1（肩を胴体重心に置く）}
$\bm a=0$ とすると $A=I_b+I_l+\mu\ell_c^2$、$B=I_l+\mu\ell_c^2$ はどちらも $q$ によらず一定になり、$\Gamma$ も一定になる。
曲率はゼロで、手足が 1 周回せればリアクションホイールと同じ働きをする（重心が回転軸からずれた「偏心ホイール」）。

\paragraph{第~\ref{ch:theory}章 2（腕の質量を 2 倍にする）}
胴体 60 kg・腕 8 kg とすると $\mu=60\times8/68=7.06$ kg、$I_l=0.327$ kg\,m$^2$。腕を真上に上げた姿勢で
$B=0.327+7.06\times0.245=2.06$、$A=5.613+0.327+7.06\times0.49=9.40$、$\Gamma=0.219$。もとの 0.142 の約 1.5 倍である（2 倍にはならない）。

\paragraph{第~\ref{ch:planar}章 1（長方形ループの置き場所）}
図~\ref{fig:D3_curvature_mass.png} の人間型で曲率が最も強いのは、腕が前方〜やや下（$-50^\circ$〜$10^\circ$）、脚が前方に上がった付近（$-20^\circ$〜$30^\circ$）である。
同じ大きさの長方形をこの付近に置き、脚の可動域の上限（30°）を超えないようにする。

\paragraph{第~\ref{ch:planar}章 2（腕の可動域を狭める）}
人間型で曲率が負の領域は腕の角度でおよそ $-110^\circ$〜$85^\circ$ に広がる。可動域を $[-90^\circ,90^\circ]$ にすると、$-110^\circ$〜$-90^\circ$ の部分が使えなくなり、上界はやや小さくなる。
正の曲率の領域（腕を後ろに回した部分）は、もともと使わないので影響しない。

\paragraph{第~\ref{ch:rw}章 1（RW と AMBAC のどちらに先に溜めるか）}
RW は周回内の揺れがなく、この人間大の例では容量も大きい（8.6 対 1.2 kg\,m$^2$/s）。姿勢の精度が重要なら RW に先に溜め、
RW が飽和に近づいたら AMBAC で追加の容量を得る（またはスラスタでアンローディングする）のが合理的である。

\paragraph{第~\ref{ch:rw}章 2（AMBAC の容量を大きくする）}
$h_{\mathrm{AMBAC}}=f_{\max}|\Delta\theta_c|/\langle1/A\rangle$ なので、(1) アクチュエータを強くして $f_{\max}$ を上げる、
(2) 手足を重く・長くする、可動域を広げる、ループの形を最適化して $|\Delta\theta_c|$ を上げる、
(3) 系全体の慣性 $A$ を大きくする（ただし $\Delta\theta_c$ は下がるので、両者の兼ね合いになる）。

\paragraph{第~\ref{ch:dyn}章 1（関節に弾性要素を入れる）}
図~\ref{fig:sim04_human_cycle.png} のパワーは周期的に正負を繰り返し、1 周の正味の仕事はゼロである。
ばねの固有振動数をループの周波数に合わせると、運動エネルギーをばねとの間でやりとりでき、アクチュエータが出す正の仕事は損失分だけに大きく減る。

\paragraph{第~\ref{ch:dyn}章 2（トルク上限が $s^2$ に比例する場合）}
$f_{\max}\propto\sqrt{\tau_{\max}/\tau_{\text{必要}}}\propto\sqrt{s^2/s^5}=s^{-3/2}$、90° 旋回時間は $s^{3/2}$ に比例する。18 m 級では人間大の約 32 倍かかる。

\paragraph{第~\ref{ch:ai}章 1（目標付近での切り替え）}
測定ノイズで誤差の符号が細かく入れ替わると、ループの回す向きがばたつく（チャタリング）。
不感帯やヒステリシスを設ける、誤差を平滑化する、目標付近では揺れのない位相で止める動作に切り替える、などの対策がある。

\paragraph{第~\ref{ch:ai}章 2（ハイブリッドの設計）}
最適制御（線形計画・MPC）が、スラスタの噴射量（特に最初の角運動量）と AMBAC の周回の予算を決める。
学習方策には、計画からのずれ（姿勢・角運動量の誤差）とループの位相を観測させ、周回周波数とスラスタの小さな補正だけを出力させる。
推進剤の節約は計画の段階で決まるので、学習方策の報酬は精度と残留回転を中心にするのがよい。

\section{第 II 部}

\paragraph{第~\ref{ch:tank}章 1（関節の速度の上限が一定の場合）}
フルード則では、大きい機体ほど動作時間が長くなる（$\sqrt{s}$ 倍）。関節の速度の上限が大きさによらず一定なら、大きい機体では上限に余裕ができ、
小さい機体ほど転倒より先に関節の速度で制限される。規格化した ZMP の曲線は一致しなくなる。

\paragraph{第~\ref{ch:tank}章 2（月面での意図しない旋回の対策）}
腕の水平の振りを 2 s 以上かけてゆっくり行う、摩擦の大きい接地面（スパイクなど）を使う、アウトリガーで接地範囲を広げる、
2 本の腕を逆向きに振って反作用を打ち消す、腕を水平面ではなく鉛直面内で動かす、などが考えられる。

\paragraph{第~\ref{ch:legs}章 1（足を前後に開く）}
足を前後に開くと、前後の支持範囲が足裏の長さより大きく広がる。静的な ZMP（約 0.13 m）が範囲に収まれば、静止して荷物を持つのに腰の戦略は不要になる。
ただし素早く運ぶときの動的な反作用には、引き続き腰の戦略が効く。

\paragraph{第~\ref{ch:legs}章 2（宇宙での脚の質量の最適値）}
脚を重くすると 1 周あたりの回転は増えるが、ループを回すのに必要なトルクも増え、$f_{\max}$ は下がる。
評価すべきは「単位時間あたりの回転」$f_{\max}|\Delta\theta_c|$ と、角運動量の容量、そして脚を重くすることで他の用途（移動・作業）に与える影響である。

\paragraph{第~\ref{ch:steps}章 1（勢いを使った段差の乗り越え）}
勢いを使う場合、必要なのは跳び上がるためのエネルギーと瞬間的なパワー、着地の衝撃に耐える強度である。
持ち上げる位置エネルギーは $mgh\propto s^4$、アクチュエータが出せるパワーは体積に比例して $s^3$ 程度なので、大型化の不利は残る。

\paragraph{第~\ref{ch:steps}章 2（前端を高くした戦車型）}
前端の高さを 0.5 m にすると、形の限界（0.25 m）より力の限界のほうが小さくなる。図~\ref{fig:D13_step_sensitivity.png} から、
摩擦が十分なら後ろへの転倒（腕を前に伸ばして 0.33 m）が限界を決め、腕による重心移動の効果がそのまま性能差になる。

\section{第 III 部}

\paragraph{第~\ref{ch:3d}章 1（3 次元での最短到達の計画）}
図~\ref{fig:D10_cone_axes.png} のような「円錐運動 1 周ごとの回転ベクトル」を動作の部品（モーションプリミティブ）として用意し、
目標の姿勢との差を、部品の回転を順に掛け合わせて埋める系列を探索する。回転が小さいうちは順序の影響は小さいので、まず大まかな配分を決め、
最後に実際の回転の合成で順序の効果を補正する方法が考えられる。

\paragraph{第~\ref{ch:3d}章 2（極を真後ろに置く）}
特異点は、腕が真後ろを向くときに生じる。背中側に手を回す動作で関節速度が発散しやすくなる。一方、頭上に手を上げる動作では特異点の問題がなくなる。
どの動作を多く使うかで、ジンバルの取り付け方向を選ぶ。

\paragraph{第~\ref{ch:shoulder}章 1（肩甲骨の動きが必要な他の姿勢）}
背中で手を組む（肩の伸展・内旋＋肩甲骨の後退）、頭の後ろで手を組む（外転・外旋＋肩甲骨の上方回旋）、
反対側の肩に手を置く（水平内転＋肩甲骨の前方突出）などがある。

\paragraph{第~\ref{ch:shoulder}章 2（モーションキャプチャーのデータを移すときの誤差）}
腕の向きだけを移すと、ひねりは経路に依存して勝手に変わるため（コッドマンのパラドックス）、手のひらや前腕の向きがずれる。
上腕の姿勢を回転（四元数など）として計測・移植し、ひねりも明示的に合わせる。測り方（基準姿勢・回転の順序）を揃えることも重要である \cite{R36}。
